% SPDX-License-Identifier: CC-BY-NC-ND-4.0
% !TEX root = main.tex

\section{圏, 関手, 自然変換}
\begin{definition} \label{def: Category}
\textbf{圏}(category)$\mathscr{C}$は以下のデータから成る.
\begin{itemize}
    \item \textbf{対象}の集まり$\Ob(\mathscr{C})$.
    \item \textbf{射}の集まり$\Mor(\mathscr{C})$.\ 射は次を満たす.
    \begin{itemize}
    \item 任意の$f \in \Mor(\mathscr{C})$は\textbf{始域}(domain)$\dom(f) \eqcolon A \in \Ob(\mathscr{C})$と\textbf{終域}(codomain)$\cod(f) \eqcolon B \in \Ob(\mathscr{C})$を持ち$f \colon A \to B$と表す.\ 始域が$A$,\ 終域が$B$であるような射の集まりを$\Mor_\mathscr{C}(A,B)$と書く.
    \item 任意の$f, g \in \Mor(\mathscr{C})$に対して$\cod(f) = \dom(g)$のとき$f$と$g$の\textbf{合成}(composition)$g \circ f \in \Mor(\mathscr{C})$が定まり$\dom(g \circ f) = \dom(f)$,\ $\cod(g \circ f) = \cod(g)$である.\ 合成は次を満たす.
        \begin{itemize}
        \item 任意の$f,g,h \in \Mor(\mathscr{C})$に対して$\cod(f) = \dom(g)$かつ$\cod(g) = \dom(h)$のとき$(h \circ g) \circ f = h \circ (g \circ f)$が成り立つ.\ この条件を\textbf{結合律}(associativity)という.
        \item 任意の$X \in \Ob(\mathscr{C})$に対して$\id_X \colon X \to X$が存在し任意の$f,g \in \Mor(\mathscr{C})$に対して$\cod(f) = \dom(g) = X$のとき$\id_X \circ f = f$かつ$g \circ \id_X = g$が成り立つ.\ $\id_X$を$X$の\textbf{恒等射}(identity)という.
        \end{itemize}
    \end{itemize}
\end{itemize}
\end{definition}

\begin{example} \label{ex: set_is_category}
対象を集合,\ 射を写像,\ 射の合成を写像の合成とすることで圏$\mathbf{Set}$が定まる.
\end{example}

\begin{definition} \label{def: OppositeCategory}
$\mathscr{C}$を圏とする.\ 対象を$\Ob(\mathscr{C})$とし$A, B \in \Ob(\mathscr{C})$に対して$\Mor_{\mathscr{C}^{\op}}(A, B) \coloneq \Mor_\mathscr{C}(B, A)$とする.\ $f \in \Mor_{\mathscr{C}^{\op}}(A, B)$と$g \in \Mor_{\mathscr{C}^{\op}}(B, C)$の$\mathscr{C}^{\op}$における合成$g \circ^{\op} f$を$\mathscr{C}$における合成$f \circ g \in \Mor_\mathscr{C}(C, A)$とすることで圏$\mathscr{C}^{\op}$が定まる.\ これを$\mathscr{C}$の\textbf{反対圏}(opposite category)という.
\end{definition}

\begin{definition} \label{def: Functor}
$\mathscr{C}, \mathscr{D}$を圏とする.\ $F$が次を満たすとき$F$は$\mathscr{C}$から$\mathscr{D}$への\textbf{共変関手}(covariant functor)であるといい$F \colon \mathscr{C} \to \mathscr{D}$と書く.
\begin{itemize}
    \item 任意の$A \in \Ob(\mathscr{C})$に対して$F(A) \in \Ob(\mathscr{D})$である.
    \item 任意の$f \in \Mor(\mathscr{C})$に対して$F(f) \in \Mor_\mathscr{D}(F(\dom(f)), F(\cod(f)))$であり以下が成り立つ.
    \begin{itemize}
    \item 任意の$X \in \Ob(\mathscr{C})$に対して$F(\id_X) = \id_{F(X)}$である.
    \item 任意の$f,g \in \Mor(\mathscr{C})$に対して$\cod(f) = \dom(g) \implies F(g \circ f) = F(g) \circ F(f)$.
    \end{itemize}
\end{itemize}
% https://q.uiver.app/#q=WzAsNSxbMCwyXSxbMCwxLCJCIl0sWzAsMCwiQSJdLFsxLDAsIkYoQSkiXSxbMSwxLCJGKEIpIl0sWzIsMSwiZiIsMl0sWzMsNCwiRihmKSJdLFs1LDYsIkYiLDAseyJzaG9ydGVuIjp7InNvdXJjZSI6MzAsInRhcmdldCI6MzB9LCJsZXZlbCI6MSwic3R5bGUiOnsidGFpbCI6eyJuYW1lIjoibWFwcyB0byJ9fX1dXQ==
\[\begin{tikzcd}
	A & {F(A)} \\
	B & {F(B)} \\
	{}
	\arrow[""{name=0, anchor=center, inner sep=0}, "f"', from=1-1, to=2-1]
	\arrow[""{name=1, anchor=center, inner sep=0}, "{F(f)}", from=1-2, to=2-2]
	\arrow["F", between={0.3}{0.7}, maps to, from=0, to=1]
\end{tikzcd}\]

\end{definition}

\begin{definition} \label{def: ContravariantFunctor}
$\mathscr{C}, \mathscr{D}$を圏とする.\ 共変関手$F \colon \mathscr{C}^{\op} \to \mathscr{D}$を\textbf{反変関手}(contravariant functor)$F \colon \mathscr{C} \to \mathscr{D}$という.
\end{definition}

\begin{definition} \label{def: constantFunctor}
$\mathscr{J}, \mathscr{C}$を圏とする.\ $A \in \Ob(\mathscr{C})$として共変関手$C_A \colon \mathscr{J} \to \mathscr{C}$を$X \in \Ob(\mathscr{J}),\ f \in \Mor(\mathscr{J})$に対して$C_A(X) \coloneq A,\ C_A(f) \coloneq \id_A$と定める.\ このとき$C_A$を\textbf{定関手}(constant functor)と呼ぶ.
\end{definition}

\begin{definition} \label{def: constantDiagramFunctor}
$\mathscr{C},\ \mathscr{J}$を圏とし共変関手$\Delta \colon \mathscr{C} \to \mathscr{C}^{\mathscr{J}}$を次のようにして定める.
\begin{itemize}
    \item $A \in \Ob(\mathscr{C})$に対して$\Delta(A) \coloneq C_A$とする.
    \item $f \in \Mor(\mathscr{C})$に対して$\Delta(f) \colon C_{\dom(f)} \Rightarrow C_{\cod(f)}$を$J \in \Ob(\mathscr{J})$に対して$\Delta(f)_J \coloneq f$で定める.
\end{itemize}
これを\textbf{定図式関手}(constant diagram functor)という.
\end{definition}

\begin{definition} \label{def: IsSmall}
$\mathscr{C}$を圏とする.\ $\Ob(\mathscr{C}) \in \Ob(\mathbf{Set})$かつ$\Mor(\mathscr{C}) \in \Ob(\mathbf{Set})$であるとき$\mathscr{C}$は\textbf{小さい圏}(small category),\ または\textbf{小圏}であるという.

任意の$A, B \in \Ob(\mathscr{C})$に対して$\Mor_\mathscr{C}(A, B) \in \Ob(\mathbf{Set})$であるとき$\mathscr{C}$は\textbf{局所的に小さい圏}(locally small category),\ または\textbf{局所小圏}であるという.
\end{definition}

\begin{example} \label{ex: hom_functor}
$\mathscr{C}$を局所小圏とし$A \in \Ob(\mathscr{C})$とする.\ $X \in \Ob(\mathscr{C})$に対して$h^A(X) \coloneq \Mor_\mathscr{C}(A, X)$とし$f \in \Mor(\mathscr{C})$に対して$h^A(f) \colon h^A(\dom(f)) \to h^A(\cod(f))$を$h^A(f)(g) \coloneq f \circ g$で定めると$h^A \colon \mathscr{C} \to \mathbf{Set}$は共変関手となる.
\end{example}

\begin{definition} \label{def: IsIsomorphism}
$\mathscr{C}$を圏とし$f \in \Mor(\mathscr{C})$とする.\ $g \in \Mor(\mathscr{C})$が存在して$\dom(f) = \cod(g)$かつ$\cod(f) = \dom(g)$となるとする.\ $f \circ g = \id_{\dom(g)}$かつ$g \circ f = \id_{\dom(f)}$となるとき$f$は\textbf{同型射}(isomorphism)であるという.\ $A,B \in \mathscr{C}$に対して同型射$f \colon A \to B$が存在するとき$A$と$B$は\textbf{同型}(isomorphic)であるといい$A \cong B$と書く.
\end{definition}

\begin{definition} \label{def: NaturalTransformation}
$F,G$を圏$\mathscr{C}$から圏$\mathscr{D}$への共変関手とする.\ $\mathscr{D}$の射の族$\eta \coloneq (\eta_X \colon F(X) \to G(X))_{X \in \Ob(\mathscr{C})}$は任意の$f \in \Mor(\mathscr{C})$に対して次の図式を可換にするとき\textbf{自然変換}(natural transformation)であるといい$\eta \colon F \Rightarrow G$と書く.\ 特に$\eta$が自然変換で任意の$X \in \Ob(\mathscr{C})$に対して$\eta_X$が同型射であるとき$\eta$は$F$から$G$への\textbf{自然同型}(natural isomorphism)であるという.\ $F$から$G$への自然変換全体の集まりを$\Nat(F, G)$と書く.

% https://q.uiver.app/#q=WzAsNCxbMCwwLCJGKFxcbWF0aHJte2RvbX0oZikpIl0sWzEsMCwiRyhcXG1hdGhybXtkb219KGYpKSJdLFswLDEsIkYoXFxtYXRocm17Y29kfShmKSkiXSxbMSwxLCJHKFxcbWF0aHJte2NvZH0oZikpIl0sWzAsMSwiXFxldGFfe1xcbWF0aHJte2RvbX0oZil9Il0sWzAsMiwiRihmKSIsMl0sWzIsMywiXFxldGFfe1xcbWF0aHJte2NvZH0oZil9Il0sWzEsMywiRyhmKSJdXQ==
\[\begin{tikzcd}
	{F(\dom(f))} & {G(\dom(f))} \\
	{F(\cod(f))} & {G(\cod(f))}
	\arrow["{\eta_{\dom(f)}}", from=1-1, to=1-2]
	\arrow["{F(f)}"', from=1-1, to=2-1]
	\arrow["{G(f)}", from=1-2, to=2-2]
	\arrow["{\eta_{\cod(f)}}", from=2-1, to=2-2]
\end{tikzcd}\]
\end{definition}

\begin{definition} \label{def: FunctorCategory}
$\mathscr{C}, \mathscr{D}$を圏とする.\ 対象を共変関手$F \colon \mathscr{C} \to \mathscr{D}$とし関手$F \colon \mathscr{C} \to \mathscr{D}$から関手$G \colon \mathscr{C} \to \mathscr{D}$への射を自然変換$\eta \colon F \Rightarrow G$とする.\ $\eta \colon F \Rightarrow G$と$\theta \colon G \Rightarrow H$の合成を$X \in \Ob(\mathscr{C})$に対して$(\theta \circ \eta)_X \coloneq \theta_X \circ \eta_X$,\ $F$の恒等射$\id_F$を$X \in \Ob(\mathscr{C})$に対して$(\id_F)_X \coloneq \id_{F(X)}$で定めることで圏$\mathscr{D}^{\mathscr{C}}$が定まる.\ これを\textbf{関手圏}(functor category)という.
\end{definition}

\begin{definition} \label{def: evaluationFunctor}
$\mathscr{C}, \mathscr{D}$を圏とする.\ $X \in \mathscr{C}$に対して共変関手$\ev_X \colon \mathscr{D}^{\mathscr{C}} \to \mathscr{D}$を$F \in \Ob(\mathscr{D}^{\mathscr{C}})$に対して$\ev_X(F) \coloneq F(X)$,\ $\eta \in \Mor_{\mathscr{D}^{\mathscr{C}}}$に対して$\ev_X(\eta) \coloneq \eta_X$で定める.\ これを\textbf{評価関手}(evaluation functor)という.
\end{definition}

% https://q.uiver.app/#q=WzAsNixbMywwLCJGKFgpIl0sWzQsMCwiRyhYKSJdLFs0LDEsIkgoWCkiXSxbMCwwLCJGIl0sWzEsMCwiRyJdLFsxLDEsIkgiXSxbMCwxLCJcXGV0YV9YIl0sWzEsMiwiXFx0aGV0YV9YIl0sWzAsMiwiKFxcdGhldGEgXFxjaXJjIFxcZXRhKV9YIiwyXSxbNCw1LCJcXHRoZXRhIiwwLHsibGV2ZWwiOjJ9XSxbMyw1LCJcXHRoZXRhIFxcY2lyYyBcXGV0YSIsMix7ImxldmVsIjoyfV0sWzMsNCwiXFxldGEiLDAseyJsZXZlbCI6Mn1dLFs5LDgsIlxcbWF0aHJte2V2fV9YIiwwLHsibGFiZWxfcG9zaXRpb24iOjQwLCJzaG9ydGVuIjp7InNvdXJjZSI6MjAsInRhcmdldCI6NDB9LCJsZXZlbCI6MSwic3R5bGUiOnsidGFpbCI6eyJuYW1lIjoibWFwcyB0byJ9fX1dXQ==
\[\begin{tikzcd}
	F & G && {F(X)} & {G(X)} \\
	& H &&& {H(X)}
	\arrow["\eta", Rightarrow, nfold, from=1-1, to=1-2]
	\arrow["{\theta \circ \eta}"', Rightarrow, nfold, from=1-1, to=2-2]
	\arrow[""{name=0, anchor=center, inner sep=0}, "\theta", Rightarrow, nfold, from=1-2, to=2-2]
	\arrow["{\eta_X}", from=1-4, to=1-5]
	\arrow[""{name=1, anchor=center, inner sep=0}, "{(\theta \circ \eta)_X}"', from=1-4, to=2-5]
	\arrow["{\theta_X}", from=1-5, to=2-5]
	\arrow["{\mathrm{ev}_X}"{pos=0.4}, between={0.2}{0.6}, maps to, from=0, to=1]
\end{tikzcd}\]
